\documentclass[UTF8,12pt,a4paper]{ctexart}

\usepackage{amsmath}
\usepackage{amssymb}
\usepackage{graphicx}
\usepackage{enumerate}
\usepackage[margin=1in]{geometry}
\geometry{a4paper}
\usepackage{booktabs}
\usepackage{array}
\usepackage{longtable}
\newcolumntype{L}[1]{>{\raggedright\arraybackslash}p{#1}}
\newcolumntype{M}[1]{>{\centering\arraybackslash}p{#1}}
\usepackage{fancyhdr}
\pagestyle{fancy}
\fancyhf{}
\usepackage{xcolor}
\usepackage{multirow}




% ------------------- 设置超链接，便于跳转 -----------------
\usepackage{enumitem}
\usepackage{hyperref}
\hypersetup{
    pdfborder={0 0 0},    % 消除超链接边框
    colorlinks=true,       % 将边框改为颜色标记（可选）
    linkcolor=black,       % 内部链接颜色（目录、引用等）
    citecolor=black,       % 引用颜色
    urlcolor=blue          % URL链接颜色（保持蓝色以便区分）
}




% ------------------------ 标题区 ----------------------------
\title{集成电路工艺技术基础 Chapter 6\\薄膜沉积详细复习手稿}
\author{冯峻}
\date{\today}
\pagenumbering{arabic}

\begin{document}



% ------------------------ 页眉和页脚 ----------------------------
\fancyhead[L]{Chapter 6 复习手稿}
\fancyhead[C]{薄膜沉积}
\fancyfoot[C]{\thepage}

\maketitle
\tableofcontents
\newpage  % \section之前用newpage分割，\subsection不用


\section{薄膜沉积概述}

\subsection{薄膜沉积的重要性}

薄膜制备是芯片加工中的一个重要工艺步骤。在现代集成电路制造中，薄膜沉积技术贯穿整个工艺流程。

\subsubsection{为什么需要薄膜沉积？}

\begin{itemize}
    \item \textbf{多层结构需求：}随着特征尺寸越来越小，需要更多的金属层来做互连
    \item \textbf{绝缘隔离：}各金属层之间的绝缘非常重要
    \item \textbf{台阶覆盖：}表面不平坦，成为VLSI多层金属高密度芯片制造的限制因素，必须考虑台阶覆盖性
    \item \textbf{可靠性：}沉积可靠的薄膜材料至关重要
\end{itemize}

\subsubsection{薄膜类型与应用}

在分立器件和集成电路制备中用到多种薄膜材料：

\begin{enumerate}
    \item \textbf{热氧化薄膜：}SiO$_2$（栅氧化层、隔离层）
    \item \textbf{电介质膜：}SiO$_2$、Si$_3$N$_4$、BPSG等
    \item \textbf{外延膜：}Si、SiGe、GaAs等
    \item \textbf{多晶硅膜：}Poly-Si（栅极、导线）
    \item \textbf{金属膜：}Al、Cu、W、Ti、Ni等（互连、接触）
\end{enumerate}

\subsection{薄膜沉积方法分类}

薄膜沉积技术主要分为两大类：\textbf{物理方法}和\textbf{化学方法}。

\subsubsection{物理方法（Physical Methods）}

\begin{itemize}
    \item \textbf{溅射（Sputtering）}
    \begin{itemize}
        \item DC溅射（Direct Current）
        \item RF溅射（Radio Frequency）
        \item 磁控溅射（Magnetron）
        \item 离子束溅射（IBS, Ion Beam Sputtering）
    \end{itemize}
    \item \textbf{蒸发（Evaporation）}
    \begin{itemize}
        \item 热蒸发（Thermal）
        \item 电子束蒸发（Electron Beam）
        \item 激光烧蚀（Laser Ablation）
    \end{itemize}
    \item \textbf{分子束外延（MBE, Molecular Beam Epitaxy）}
    \item \textbf{离子镀（Ion Plating）}
\end{itemize}

\subsubsection{化学方法（Chemical Methods）}

\begin{itemize}
    \item \textbf{化学气相沉积（CVD, Chemical Vapor Deposition）}
    \begin{itemize}
        \item 常压CVD（APCVD）
        \item 低压CVD（LPCVD）
        \item 等离子体增强CVD（PECVD）
        \item 金属有机CVD（MOCVD）
    \end{itemize}
    \item \textbf{原子层沉积（ALD, Atomic Layer Deposition）}
    \item \textbf{电镀（Electroplating）}与无电镀（Electroless）
    \item \textbf{溶胶-凝胶法（Sol-gel）}
    \item \textbf{旋涂（Spin Coating）}
    \item \textbf{浸渍（Dipping）}
    \item \textbf{喷涂（Spraying）}
\end{itemize}

\subsection{薄膜的关键性能指标}

\subsubsection{薄膜性能}

\begin{itemize}
    \item \textbf{应力（Stress）：}内应力会导致薄膜开裂或剥离
    \item \textbf{附着性（Adhesion）：}薄膜与衬底之间的粘附强度
    \item \textbf{化学计量比（Stoichiometry）：}化合物薄膜的组分比例
    \item \textbf{密度（Density）：}影响薄膜的机械和电学性能
    \item \textbf{晶粒尺寸、晶界特性、取向（Grain size, boundary properties, orientation）}
    \item \textbf{电学性能：}如击穿电压等
    \item \textbf{光学性能：}折射率、透过率等
    \item \textbf{磁学性能}
\end{itemize}

\subsubsection{工艺要求}

\begin{itemize}
    \item \textbf{成分控制：}所需的组分比例
    \item \textbf{低污染：}避免杂质引入
    \item \textbf{均匀性：}
    \begin{itemize}
        \item 片内均匀性（thickness/morphology across sample）
        \item 片间均匀性（thickness/morphology run to run）
    \end{itemize}
    \item \textbf{方向性：}
    \begin{itemize}
        \item 定向沉积（directional）：有利于lift-off工艺和沟槽填充
        \item 非定向沉积（non-directional）：有利于台阶覆盖
    \end{itemize}
\end{itemize}

\section{物理气相沉积（PVD）}

\subsection{PVD基本原理}

\textbf{物理气相沉积（PVD, Physical Vapor Deposition）}是通过物理手段将材料转化为气相，然后在衬底上凝结形成薄膜的技术。

\subsubsection{PVD的基本步骤}

所有PVD工艺都按照以下顺序进行：

\begin{enumerate}
    \item \textbf{气化：}待沉积材料通过物理手段转化为气相
    \item \textbf{输运：}蒸气穿过低压区域（从源到衬底）
    \item \textbf{凝结：}蒸气在衬底上凝结形成薄膜
\end{enumerate}

\subsubsection{PVD系统要素}

\begin{itemize}
    \item \textbf{真空系统（Vacuum）：}提供低压环境，减少气体分子碰撞
    \item \textbf{靶材/蒸发源：}提供沉积材料
    \item \textbf{衬底及加热/冷却系统：}控制薄膜生长条件
    \item \textbf{等离子体系统（针对溅射）：}产生离子轰击靶材
\end{itemize}

\section{溅射技术（Sputtering）}

\subsection{溅射基本原理}

\subsubsection{什么是溅射？}

\textbf{溅射（Sputtering）}是一种通过高能粒子轰击靶材表面，使靶材原子被打出并沉积到衬底上的工艺。

\textbf{基本过程：}
\begin{enumerate}
    \item 在真空腔体中产生等离子体，生成高能离子（通常为Ar$^+$）
    \item 离子被电场加速并轰击靶材（源材料）
    \item 靶材表面原子被撞击后脱离，称为"溅射"
    \item 被溅射的原子输运到衬底表面
    \item 原子在衬底上凝结形成薄膜
\end{enumerate}

\subsubsection{溅射阈值与溅射产额}

\begin{itemize}
    \item \textbf{溅射阈值（Sputtering Threshold）}
    \begin{itemize}
        \item 能量范围：10-30 eV
        \item 取决于靶材材料种类
        \item 低于此能量的离子无法引起溅射
    \end{itemize}
    
    \item \textbf{溅射产额/溅射率（Sputtering Yield）}
    \begin{itemize}
        \item 定义：平均每个入射离子能打出的靶材原子数
        \item 影响因素：
        \begin{enumerate}
            \item 入射离子能量
            \item 入射离子种类
            \item 入射角度
            \item 靶材材料
        \end{enumerate}
        \item 特点：溅射率随靶材原子序数呈周期性变化
    \end{itemize}
\end{itemize}

\subsection{气体辉光放电}

溅射过程需要通过辉光放电产生等离子体。典型的直流辉光放电曲线可分为多个区间：

\begin{itemize}
    \item \textbf{DE段（辉光放电区）：}自持放电段，电流增加时放电电压低且保持不变
    \item \textbf{EF段（反常辉光放电段/溅射区）：}电流增加时放电电压升高，辉光已占据整个阴极。这是溅射工艺的工作区间。
\end{itemize}

\subsection{溅射分类}

\subsubsection{直流溅射（DC Sputtering）}

\textbf{工作原理：}
\begin{enumerate}
    \item 自由电子与Ar原子碰撞
    \begin{itemize}
        \item 激发Ar原子
        \item 电离Ar原子产生Ar$^+$和二次电子
    \end{itemize}
    \item 二次电子产生更多等离子体
    \item Ar$^+$被加速向靶材运动，溅射靶材
    \item 沉积过程是非定向的（除非使用准直器）
\end{enumerate}

\textbf{特点：}
\begin{itemize}
    \item 结构简单
    \item \textbf{只能用于导电靶材}（金属等）
    \item 沉积速率较低
\end{itemize}

\subsubsection{射频溅射（RF Sputtering）}

\textbf{射频电场放电的特点：}
\begin{itemize}
    \item 电子在两极间振荡运动，使气体容易电离
    \item 放电电压低
    \item \textbf{电极可以是导体或绝缘体}
\end{itemize}

\textbf{工作原理：}
\begin{itemize}
    \item RF电场中和电极上的正电荷积累
    \item \textbf{可以沉积介电材料}（绝缘体）
    \item 在低频（< 100 kHz）时，电子和离子都能跟随电场极性变化
    \item 在高频（> 1 MHz）时，重离子无法跟随电场切换
\end{itemize}

\textbf{典型工作频率：}13.56 MHz（工业标准）

\subsubsection{磁控溅射（Magnetron Sputtering）}

\textbf{工作原理：}
\begin{itemize}
    \item 磁场将电子限制在靶材区域附近
    \item 电子在磁场中做螺旋运动，增加与气体分子的碰撞概率
\end{itemize}

\textbf{优点：}
\begin{itemize}
    \item 电流密度增加
    \item 放电压力降低（典型值：1-20 mTorr）
    \item 沉积速率提高
\end{itemize}

\textbf{关键挑战：}设计磁场分布以获得靶材全表面的均匀刻蚀。通过合理的磁铁排列，可以控制刻蚀形貌并获得良好的沉积均匀性。

\subsection{溅射薄膜的成核与生长}

\subsubsection{薄膜生长模式}

薄膜生长模式取决于衬底与薄膜之间的相互作用：

\begin{enumerate}
    \item \textbf{层状生长（Frank-van der Merwe模式）}
    \begin{itemize}
        \item 条件：衬底与薄膜晶格匹配良好，$\mu_{sf} > \mu_{ss}$
        \item 特点：逐层生长，形成连续薄膜
    \end{itemize}
    
    \item \textbf{岛状生长（Volmer-Weber模式）}
    \begin{itemize}
        \item 条件：衬底与薄膜晶格失配大，$\mu_{sf} < \mu_{ss}$
        \item 特点：形成三维岛状结构
    \end{itemize}
    
    \item \textbf{层加岛状生长（Stranski-Krastanov模式）}
    \begin{itemize}
        \item 介于(a)和(c)之间
        \item 先形成几层润湿层，然后转为岛状生长
    \end{itemize}
\end{enumerate}

\subsubsection{柱状晶生长}

多晶薄膜中常见\textbf{柱状晶生长（Columnar Grain Growth）}现象，晶粒呈柱状垂直于衬底生长。

\subsubsection{结构区域图（Zone Diagram）}

薄膜微观结构与衬底温度密切相关。使用MD（Movchan-Demchishin）区域图可以预测薄膜结构：

\begin{itemize}
    \item \textbf{Zone 1：}低温区，原子表面迁移受限，形成多孔柱状结构
    \item \textbf{Zone T（过渡区）：}纤维状结构
    \item \textbf{Zone 2：}中温区，柱状晶结构，晶界清晰
    \item \textbf{Zone 3：}高温区，体扩散显著，等轴晶结构
\end{itemize}

\section{蒸发技术（Evaporation）}

\subsection{蒸发基本原理}

\textbf{蒸发（Evaporation）}是通过加热使源材料蒸发，蒸气原子/分子在真空中输运并在衬底上凝结成膜的技术。

\subsection{蒸发方式分类}

\subsubsection{热蒸发（Thermal Evaporation）}

\begin{itemize}
    \item 使用电阻加热使材料蒸发
    \item 适用于低熔点材料
    \item 结构简单，成本低
\end{itemize}

\subsubsection{电子束蒸发（E-beam Evaporation）}

\begin{itemize}
    \item 使用高能电子束加热材料
    \item \textbf{适用于高熔点材料}（Ti、W等）
    \item 沉积材料：Ti、Al、Cu、TiN、TaN等
    \item 配备石英晶体微天平（Quartz Crystal Microbalance）监测沉积速率
\end{itemize}

\subsubsection{其他加热方式}

\begin{itemize}
    \item 激光加热
    \item 高频感应加热
\end{itemize}

\subsection{多组分蒸发}

对于合金或化合物薄膜的蒸发，存在以下挑战：

\begin{itemize}
    \item 各组分材料的蒸气压不同
    \item 蒸发速率不同导致薄膜组分偏离
\end{itemize}

\textbf{解决方案：}
\begin{enumerate}
    \item 选择各组分蒸气压接近的合金靶
    \item 蒸发结束后进行高温退火
    \item 不同温度下多源同时蒸发（共蒸发）
\end{enumerate}

\subsection{热蒸发与电子束蒸发比较}

\begin{longtable}{|L{3.5cm}|L{5.5cm}|L{5.5cm}|}
\hline
\textbf{特性} & \textbf{热蒸发} & \textbf{电子束蒸发} \\
\hline
\endfirsthead
\hline
\textbf{特性} & \textbf{热蒸发} & \textbf{电子束蒸发} \\
\hline
\endhead
\hline
\endfoot
\hline
\endlastfoot

加热方式 & 电阻加热 & 电子束轰击 \\
\hline
适用材料 & 低熔点材料 & 高熔点材料 \\
\hline
能量效率 & 较低 & 较高 \\
\hline
污染 & 可能有坩埚污染 & 局部加热，污染小 \\
\hline
设备成本 & 低 & 高 \\
\hline
沉积速率 & 中等 & 较高 \\
\hline
\end{longtable}

\subsection{蒸发与溅射比较}

\begin{longtable}{|L{3cm}|L{5.5cm}|L{5.5cm}|}
\hline
\textbf{特性} & \textbf{蒸发} & \textbf{溅射} \\
\hline
\endfirsthead
\hline
\textbf{特性} & \textbf{蒸发} & \textbf{溅射} \\
\hline
\endhead
\hline
\endfoot
\hline
\endlastfoot

原理 & 热能使材料蒸发 & 离子轰击使原子溅出 \\
\hline
沉积方向性 & 高（直线运动） & 可调（准直器控制） \\
\hline
台阶覆盖 & 较差 & 较好 \\
\hline
材料选择 & 受蒸气压限制 & 几乎所有材料 \\
\hline
薄膜附着力 & 较弱 & 较强 \\
\hline
合金沉积 & 组分控制困难 & 组分保持较好 \\
\hline
沉积温度 & 较高 & 可以较低 \\
\hline
\end{longtable}

\section{工艺控制（Process Control）}

\subsection{主要工艺参数}

PVD工艺的主要控制参数包括：

\begin{enumerate}
    \item \textbf{压力（Pressure）}
    \item \textbf{温度（Temperature）}
    \item \textbf{功率（Power）}
    \item \textbf{间距（Spacing）}——靶材与衬底之间的距离
\end{enumerate}

\subsection{温度控制}

\subsubsection{衬底温度控制的重要性}

\begin{itemize}
    \item 溅射过程中必须精确控制晶圆温度
    \item 温度影响薄膜的晶粒尺寸和微观结构
\end{itemize}

\subsubsection{温度控制方法}

最常用的方法是\textbf{背面热传导}：
\begin{itemize}
    \item 在晶圆背面施加高压Ar气
    \item 通过导热传输实现晶圆与加热器之间的热耦合
    \item 加热器位于晶圆下方
\end{itemize}

\subsubsection{温度对薄膜的影响}

\begin{itemize}
    \item 衬底温度升高，薄膜晶粒尺寸增大
    \item 表面原子迁移率增加
    \item 影响薄膜致密度和应力
\end{itemize}

\subsection{表面清洁度的影响}

衬底表面清洁度对薄膜晶粒尺寸有显著影响：
\begin{itemize}
    \item 清洁表面有利于均匀成核
    \item 污染物可能成为非均匀形核点
    \item 影响薄膜与衬底的附着力
\end{itemize}

\subsection{压力控制与测量}

\subsubsection{压力单位换算}

\begin{itemize}
    \item 1 atm = 760 Torr = 1 Bar = $10^5$ Pa = 14.7 psi
    \item 1 mbar = 160 mTorr
    \item 1 Pa = 7.6 mTorr
\end{itemize}

\subsubsection{压力测量仪器}

\begin{enumerate}
    \item \textbf{对流真空计（Convection Gauge）}
    \begin{itemize}
        \item 测量范围：大气压 - 50 mTorr
        \item 用于粗真空测量
    \end{itemize}
    
    \item \textbf{薄膜真空计（Baratron）}
    \begin{itemize}
        \item 测量范围：100 mTorr - 1 mTorr
        \item 用于工艺压力测量
    \end{itemize}
    
    \item \textbf{电离真空计（Ionization Gauge）}
    \begin{itemize}
        \item 测量范围：100 mTorr - $10^{-9}$ Torr
        \item 用于基底压力测量
        \item 分为热阴极和冷阴极两种
    \end{itemize}
\end{enumerate}

\subsection{台阶覆盖（Step Coverage）}

\subsubsection{台阶覆盖的重要性}

在沟槽（trench）或通孔（via）中沉积薄膜时，台阶覆盖是关键挑战：
\begin{itemize}
    \item 随着特征尺寸缩小，深宽比（aspect ratio）增加
    \item 侧壁和底部的沉积均匀性难以保证
    \item 可能出现空洞（void）和悬挂（overhang）
\end{itemize}

\subsubsection{台阶覆盖的评价指标}

\begin{itemize}
    \item 侧壁覆盖率
    \item 底部覆盖率
    \item 填充完整性
\end{itemize}

\subsubsection{改善台阶覆盖的方法}

\textbf{准直溅射（Collimated Sputtering）：}
\begin{itemize}
    \item 在靶材和衬底之间放置准直器（Collimator）
    \item 过滤掉大角度入射的原子
    \item 提高沉积的方向性
    \item 改善深沟槽和通孔的底部覆盖
\end{itemize}

\section{PVD应用}

PVD技术在集成电路制造中有广泛应用：

\subsection{主要应用领域}

\begin{enumerate}
    \item \textbf{自对准硅化物（Salicide）}
    \begin{itemize}
        \item 沉积Ti、Co、Ni等金属
        \item 与Si反应形成低阻硅化物
        \item 降低源/漏和栅极电阻
    \end{itemize}
    
    \item \textbf{阻挡层（Barrier）}
    \begin{itemize}
        \item Ti/TiN、Ta/TaN等
        \item 防止Cu向介电层扩散
        \item 改善附着性
    \end{itemize}
    
    \item \textbf{抗反射涂层（ARC, Anti-Reflective Coating）}
    \begin{itemize}
        \item TiN等材料
        \item 减少光刻过程中的反射
        \item 改善图形精度
    \end{itemize}
    
    \item \textbf{互连（Interconnect）}
    \begin{itemize}
        \item Al、Cu金属布线
        \item 多层金属互连
    \end{itemize}
\end{enumerate}

\section{章节总结与知识点梳理}

\subsection{核心概念对照}

\begin{longtable}{|L{3.5cm}|L{11cm}|}
\hline
\textbf{概念} & \textbf{关键要点} \\
\hline
\endfirsthead
\hline
\textbf{概念} & \textbf{关键要点} \\
\hline
\endhead
\hline
\endfoot
\hline
\endlastfoot

薄膜沉积分类 & 物理方法（PVD）和化学方法（CVD等） \\
\hline

PVD基本步骤 & 气化 $\rightarrow$ 输运 $\rightarrow$ 凝结 \\
\hline

溅射原理 & 高能离子轰击靶材，打出原子沉积到衬底 \\
\hline

溅射阈值 & 10-30 eV，取决于靶材材料 \\
\hline

溅射产额 & 每个入射离子打出的原子数，与能量、入射角、材料有关 \\
\hline

DC溅射 & 只能用于导电材料，结构简单 \\
\hline

RF溅射 & 可沉积绝缘材料，频率13.56 MHz \\
\hline

磁控溅射 & 磁场限制电子，提高电离率和沉积速率 \\
\hline

薄膜生长模式 & 层状、岛状、层加岛状，取决于界面能 \\
\hline

热蒸发 & 电阻加热，适用于低熔点材料 \\
\hline

电子束蒸发 & 高能电子束加热，适用于高熔点材料 \\
\hline

PVD工艺参数 & 压力、温度、功率、间距 \\
\hline

台阶覆盖 & 沟槽/通孔沉积的关键挑战 \\
\hline

准直溅射 & 使用准直器提高方向性，改善底部覆盖 \\
\hline

PVD应用 & Salicide、阻挡层、ARC、互连 \\
\hline
\end{longtable}

\subsection{重要数据和参数}

\begin{enumerate}
    \item 溅射阈值能量：10-30 eV
    \item RF溅射标准频率：13.56 MHz
    \item 磁控溅射典型压力：1-20 mTorr
    \item 压力换算：1 atm = 760 Torr = $10^5$ Pa
    \item 1 mbar = 160 mTorr；1 Pa = 7.6 mTorr
\end{enumerate}

\subsection{常见考点和易错点}

\begin{enumerate}
    \item \textbf{DC溅射与RF溅射的区别}
    \begin{itemize}
        \item DC：只能用于导体
        \item RF：可用于绝缘体，通过交变电场中和电荷积累
    \end{itemize}
    
    \item \textbf{磁控溅射的优势}
    \begin{itemize}
        \item 提高电离效率
        \item 降低工作压力
        \item 提高沉积速率
        \item 减少衬底损伤
    \end{itemize}
    
    \item \textbf{蒸发与溅射的选择}
    \begin{itemize}
        \item 蒸发：方向性好，适合lift-off
        \item 溅射：台阶覆盖好，附着力强
    \end{itemize}
    
    \item \textbf{温度对薄膜的影响}
    \begin{itemize}
        \item 温度升高，晶粒尺寸增大
        \item 影响薄膜密度和应力
        \item 根据区域图预测薄膜结构
    \end{itemize}
    
    \item \textbf{台阶覆盖问题}
    \begin{itemize}
        \item 高深宽比结构难以填充
        \item 准直器可改善底部覆盖
    \end{itemize}
\end{enumerate}

\subsection{与其他章节的联系}

\begin{itemize}
    \item \textbf{与清洗工艺：}沉积前需要清洁的衬底表面确保粘附性
    \item \textbf{与光刻工艺：}PVD金属层需要光刻定义图形
    \item \textbf{与刻蚀工艺：}沉积后通过刻蚀形成所需图形
    \item \textbf{与氧化工艺：}热氧化是另一种形成SiO$_2$的方法
    \item \textbf{与CVD：}CVD是另一种重要的薄膜沉积方法，台阶覆盖更好
\end{itemize}

\section{复习建议}

\subsection{重点掌握内容}

\begin{enumerate}
    \item 薄膜沉积方法的分类（物理方法vs化学方法）
    \item PVD的三个基本步骤
    \item 溅射的基本原理和过程
    \item DC溅射、RF溅射、磁控溅射的原理和特点
    \item 薄膜生长的三种模式
    \item 热蒸发与电子束蒸发的比较
    \item 蒸发与溅射的比较
    \item 主要工艺参数及其影响
    \item PVD的主要应用
\end{enumerate}

\subsection{理解性内容}

\begin{enumerate}
    \item 为什么DC溅射不能用于绝缘材料？
    \item 为什么RF溅射可以沉积绝缘材料？
    \item 磁场如何提高溅射效率？
    \item 衬底温度如何影响薄膜微观结构？
    \item 为什么需要准直器来改善台阶覆盖？
    \item 溅射产额与哪些因素有关？
\end{enumerate}

\subsection{记忆技巧}

\begin{itemize}
    \item \textbf{PVD步骤：}气化$\rightarrow$输运$\rightarrow$凝结（三步法）
    \item \textbf{溅射类型：}DC（导体）$\rightarrow$ RF（绝缘体）$\rightarrow$ 磁控（高效率）
    \item \textbf{生长模式：}层状（匹配好）$\rightarrow$ 岛状（失配大）$\rightarrow$ 层加岛（过渡）
    \item \textbf{蒸发vs溅射：}蒸发方向性好，溅射覆盖好
\end{itemize}

\end{document}
